% R. Brown's translations of two letters to L. Breen, with his note and comments (cote 134-2): archive cote 134-2, site page 139 (PDF page 140).
% translation of the letter of 17-19 July 1975. Draft compared: ps1.tex, lines 1605-1614.
% AI-transcribed by Claude Opus 5.5 (claude-opus-5-5, adaptive thinking, effort high), Message Batches API, batch brown_full (msgbatch_01Md3i5yK1PxkpJn5GHfEExE), 2026-11-13; not yet checked by eye.
\scanpage{134-2/139}
\mspage{- 12 -}

and with verifying the known formulas for these operations).
But it appears probable that, to carry out this transcription well,
it would be suitable, to regard a diagram of type
\diagram{\begin{tikzcd}
  \cdot \ar[r]\ar[d] & {} \\
  {} &
\end{tikzcd}}
as based in the homotopy category formed from the category of
semi-simplicial sets, i.e. to consider the category of diagrams
of semi-simplicial sets, and to form from this a homotopy category
of functions.

I have recently more or less made explicit, while thinking
on \unsure{tne|the} foundations of ``tame topology'',(i.e. where one immediately
eliminates all wild phenomena) how an equisingular stratification,
say with non singular strata, of a compact ``tame  space'', gives
rise canonically to a diagram of spaces which are manifolds with
boundary, the arrows of the diagram being essentially locally
trivial fibrations of manifolds with boundary on the others (with
fibres manifolds with boundary), \uline{and} the inclusion of the boundary
in these manifolds with boundary [in fact, one finds slightly
more general inclusions, certain boundaries which appear being
provided with ``elementary'' cellular decomposition, i.e. the
closed strata are like manifolds with boundary which paste along
some of the common part of the boundary; and there is
% NOTE: the typist left the left half of the next line blank (a gap in the typed text, nothing hidden)
\hspace*{20em} the inclusions of these pieces
one in another ...], and can be reconstituted from this diagram
by pasting.  In other words, one has a canonical devissage
description of a tame compact space  $X$ , eventually union of an
equisingular stratification with non-singular strata, in terms
of finite diagrams of a precise nature made out of manifolds

\mspage{\hand{140}}
